\medskip

Pour chacune des affirmations suivantes, indiquer si elle est vraie ou fausse. 

Chaque réponse doit être justifiée. 

Une réponse non justifiée ne rapporte aucun point.

\medskip

\begin{enumerate}
\item \textbf{Affirmation :} La fonction $f$ définie sur $\R$ par $f(x) = \text{e}^x - x$ est convexe.
\item \textbf{Affirmation :} L'équation $\left(2\text{e}^x - 6\right)\left(\text{e}^x + 2\right) = 0$ admet $\ln (3)$ comme unique solution dans $\R$.
\item \textbf{Affirmation :}
\[\displaystyle\lim_{x \to + \infty} \dfrac{\text{e}^{2x} - 1}{\text{e}^x - x} = 0.\]

\item Soit $f$ la fonction définie sur $\R$ par $f(x) = (6x + 5)\text{e}^{3x}$ et $F$ la fonction définie sur $\R$ par : 

$F(x) = (2x + 1)\text{e}^{3x} + 4$.

\textbf{Affirmation :} $F$ est la primitive de $f$ sur $\R$ qui prend la valeur 5 quand $x = 0$.
\item On considère la fonction \texttt{mystere} définie ci-dessous qui prend une liste L de nombres en paramètre.

On rappelle que \texttt{len}(L) représente la longueur de la liste L.

\begin{center}
\begin{tabular}{|l|}\hline
\textbf{def mystere(L)} :\\
	\quad  S $= 0$\\
	\quad \textbf{for} i in range(\texttt{len}(L)) :\\
	\quad \qquad S = S + L[i]\\
	\quad \textbf{return} S / \texttt{len}(L)\\ \hline
	\end{tabular}
\end{center}

\textbf{Affirmation :} L'exécution de \texttt{mystere}([1, 9, 9, 5, 0, 3, 6,  12, 0, 5]) renvoie $50$.
\end{enumerate}

\bigskip
